\newpage
\section{\textsc{Analysis and Design}}
\label{chap:arch}
\rhead{\itshape Chapter 2: Analysis and Design}
\lhead{}
\cfoot{\thepage}
\pagestyle{fancy}
\renewcommand{\headrulewidth}{0.4pt}

\begin{flushright}
\textit{``Every great design begins with an even better story.''\\
Lorinda Mamo}
\end{flushright}

\subsection{Project Genesis: Iterative Development}

ExamGen was not designed in its final form. It evolved through two prior versions,
each revealing specific limitations that were later addressed by findings in the
reference literature.

\subsubsection{Version 1 \& 2 --- From Stubs to Failed Prototype}

V1 established the school management platform (auth, courses, PDF uploads) with
AI stubs. V2 introduced a RAG-Phi4 generation component that revealed three
failures, each driving a corrective decision in V3:

Three failures were identified and each drove a specific corrective decision:

\begin{description}
  \item[F1 --- No structured output.] Free-text paragraphs instead of JSON;
  regex parsing failed on $\approx$30\% of responses.
  \textit{Fix:} wrap every LLM call with \texttt{instructor} \cite{instructor2023}
  (self-healing Pydantic schema enforcement, \texttt{max\_retries=3}).

  \item[F2 --- Free-tier model fidelity.] Models ignored structural constraints
  (e.g.\ ``exactly 4 choices'').
  \textit{Fix:} switch to paid-tier DeepSeek-V4-Flash \cite{deepseek2024}
  with fallback DeepSeek $\to$ Gemini $\to$ Claude.

  \item[F3 --- Semantic retrieval blends topics.] Retrieved passages mixed subjects
  and years, diluting the instructor's style.
  \textit{Fix:} separate style anchors (filter-based sampling: same teacher, subject,
  Bloom level) from RAG factual context --- two distinct prompt blocks.
\end{description}

\subsection{System Sequence Diagram}
\label{subsec:sequence}

Figure~\ref{fig:sysseq} traces the complete runtime interaction across the five
system components for the two main workflows: (1)~PDF upload and parsing
through the five-stage \texttt{docParser} pipeline, and (2)~exam generation
with hybrid RAG retrieval, LLM generation, and post-generation evaluation.

\begin{figure}[H]
\centering
\resizebox{\textwidth}{!}{%
\begin{sequencediagram}

\newthread[white]{t}{\textbf{Teacher}}
\newinst[1.6]{fe}{\textbf{Frontend}}
\newinst[1.6]{be}{\textbf{Backend}}
\newinst[1.6]{dp}{\textbf{docParser}}
\newinst[1.6]{ef}{\textbf{exam-forge}}

%% ── Phase 1: Authentication ──────────────────────────────────────────
\begin{call}{t}{login / register}{fe}{JWT session}
  \begin{call}{fe}{POST /auth/login}{be}{JWT token}
  \end{call}
\end{call}

\postlevel

%% ── Phase 2: Document Upload and Parsing ─────────────────────────────
\begin{call}{t}{uploadPDF(file, subject)}{fe}{status: processing}
  \begin{call}{fe}{POST /courses/upload}{be}{course saved}
    \begin{call}{be}{parseAsync(path) [background]}{dp}{JSON + confidence c}
      \begin{call}{dp}{S1: Docling to Markdown + PyMuPDF fill}{dp}{}
      \end{call}
      \begin{call}{dp}{S2: LLM + instructor to Exam JSON}{dp}{}
      \end{call}
      \begin{call}{dp}{S3: Grounding validate + confidence c}{dp}{}
      \end{call}
      \begin{call}{dp}{S4: Bloom classify (GPT-OSS-120B, temp=0)}{dp}{}
      \end{call}
      \begin{call}{dp}{S5: VLM diagram transcription (Qwen2.5-VL)}{dp}{}
      \end{call}
    \end{call}
    \begin{call}{be}{POST /api/index-file}{ef}{RAG indexed}
    \end{call}
  \end{call}
\end{call}

\postlevel

%% ── Phase 3: Exam Generation ─────────────────────────────────────────
\begin{call}{t}{generateExam(subj, level, nq)}{fe}{SSE stream}
  \begin{call}{fe}{POST /api/generate}{be}{}
    \begin{call}{be}{planExam()}{ef}{ExamPlan}
      \begin{call}{ef}{TemplateAnalyzer: mine archive patterns}{ef}{}
      \end{call}
    \end{call}
    \begin{call}{be}{assembleExam(plan) [loop per question]}{ef}{SSE events}
      \begin{call}{ef}{HybridRetrieve: BM25 + FAISS + RRF}{ef}{}
      \end{call}
      \begin{call}{ef}{sampleStyleAnchors(subj, level)}{ef}{}
      \end{call}
      \begin{call}{ef}{generate(ctx, style): DeepSeek-V4-Flash}{ef}{}
      \end{call}
      \begin{call}{ef}{evaluate(question): GPT-OSS 4-axis}{ef}{}
      \end{call}
    \end{call}
    \mess{be}{SSE: question + eval result}{fe}
  \end{call}
\end{call}

\postlevel

%% ── Phase 4: Review and Export ───────────────────────────────────────
\mess{fe}{card: accept / edit / reject}{t}
\begin{call}{t}{exportExam()}{fe}{PDF archived}
  \begin{call}{fe}{POST /api/archive}{be}{saved}
  \end{call}
\end{call}

\end{sequencediagram}
}%
\caption{UML Sequence Diagram of the complete ExamGen workflow.
Phase~1: authentication. Phase~2: PDF upload and five-stage IDP parsing.
Phase~3: hybrid RAG retrieval, LLM question generation, and 4-axis evaluation
with SSE streaming. Phase~4: instructor review and exam archiving.}
\label{fig:sysseq}
\end{figure}

\subsection{Final System Architecture}

ExamGen comprises four layers communicating through HTTP and filesystem interfaces,
as illustrated in Figure~\ref{fig:arch}.

\begin{figure}[H]
\centering
\scalebox{0.82}{%
\begin{tikzpicture}[
  font=\small,
  % Component box style
  layer/.style={
    draw=black!80, thick, rounded corners=4pt,
    minimum width=11.5cm, align=center,
    inner sep=8pt
  },
  % Title bar inside each layer
  ltitle/.style={font=\small\bfseries},
  % Connector arrow
  conn/.style={->, >=latex, thick, gray!70},
  % Small sub-component inside a layer
  sub/.style={
    draw=black!50, rounded corners=2pt,
    fill=white, inner sep=4pt,
    font=\scriptsize, align=center,
    minimum width=2cm, minimum height=0.6cm
  }
]

%% ── LAYER 1 ────────────────────────────────────────────────────────────────
\node[layer, fill=gray!12] (L1) at (0,0) {
  \begin{tikzpicture}[font=\scriptsize, inner sep=2pt]
    \node[ltitle] {\textsf{Layer 1 — Raw Document Corpus}\ (\texttt{examQ/})};
    \node[below=0.15cm, font=\scriptsize]
      {137 files: PDF · DOCX · JPEG \quad|\quad 5 subjects \quad|\quad French \& Arabic};
  \end{tikzpicture}
};

%% ── ARROW 1→2 ──────────────────────────────────────────────────────────────
\node[below=0.05cm of L1, font=\scriptsize\itshape, text=gray]
  (arr12) {PDF upload via web $\;\to\;$ \texttt{docparser\_service} launches background parse};
\draw[conn] (L1.south) -- (arr12.north);

%% ── LAYER 2 ────────────────────────────────────────────────────────────────
\node[layer, fill=blue!8, below=0.55cm of L1] (L2) {
  \begin{tikzpicture}[font=\scriptsize, node distance=0.25cm and 0.3cm]
    \node[ltitle] (t2) {\textsf{Layer 2 — Intelligent Document Processing}\ (\texttt{docParser})};
    \node[sub, fill=blue!15, below=0.2cm of t2]        (s1) {S1: Docling\\PDF$\to$Markdown};
    \node[sub, fill=blue!15, right=0.25cm of s1]       (s2) {S2: LLM+Instructor\\$\to$ Exam JSON};
    \node[sub, fill=blue!15, right=0.25cm of s2]       (s3) {S3: Grounding\\Validator};
    \node[sub, fill=blue!15, right=0.25cm of s3]       (s4) {S4: Bloom\\Classifier};
    \node[sub, fill=blue!15, right=0.25cm of s4]       (s5) {S5: VLM\\Diagrams};
    \draw[->, >=latex, thick, blue!60]
      (s1) -- (s2) -- (s3) -- (s4) -- (s5);
    \node[below=0.2cm of s3, font=\scriptsize\itshape]
      {Output $\to$ \texttt{exam-forge/data/final\_unified/teacher\_\{id\}/}};
  \end{tikzpicture}
};

%% ── ARROW 2→3 ──────────────────────────────────────────────────────────────
\node[below=0.05cm of L2, font=\scriptsize\itshape, text=gray]
  (arr23) {HTTP:\ \texttt{/api/index-file}\ (cloud RAG) \;+\; \texttt{/api/invalidate-cache}};
\draw[conn] (L2.south) -- (arr23.north);

%% ── LAYER 3 ────────────────────────────────────────────────────────────────
\node[layer, fill=green!7, below=0.55cm of L2] (L3) {
  \begin{tikzpicture}[font=\scriptsize, node distance=0.25cm and 0.35cm]
    \node[ltitle] (t3)
      {\textsf{Layer 3 — RAG-Augmented Generation Engine}\ (\texttt{exam-forge})};
    \node[sub, fill=green!20, below=0.2cm of t3]         (r) {Retrieval\\BM25+AraBERT/MiniLM\\RRF fusion};
    \node[sub, fill=green!20, right=0.35cm of r]          (g) {Generator\\DeepSeek-V4-Flash\\MCQ/SAQ/Exercise};
    \node[sub, fill=green!20, right=0.35cm of g]          (e) {Evaluator\\4-axis quality\\check};
    \node[sub, fill=green!20, right=0.35cm of e]          (p) {Planner\\TemplateAnalyzer\\ExamAssembler};
    \draw[->, >=latex, thick, green!60!black]
      (r) -- (g) -- (e) -- (p);
    \node[below=0.2cm of g, font=\scriptsize\itshape]
      {FastAPI · Server-Sent Events (SSE) streaming};
  \end{tikzpicture}
};

%% ── ARROW 3→4 ──────────────────────────────────────────────────────────────
\node[below=0.05cm of L3, font=\scriptsize\itshape, text=gray]
  (arr34) {REST API / SSE};
\draw[conn] (L3.south) -- (arr34.north);

%% ── LAYER 4 ────────────────────────────────────────────────────────────────
\node[layer, fill=orange!8, below=0.55cm of L3] (L4) {
  \begin{tikzpicture}[font=\scriptsize, node distance=0.25cm and 0.5cm]
    \node[ltitle] (t4)
      {\textsf{Layer 4 — Web Application}\ (\texttt{app/project/})};
    \node[sub, fill=orange!20, below=0.2cm of t4]        (b) {Backend\\FastAPI + SQLite\\Teacher/Subject/Course};
    \node[sub, fill=orange!20, right=0.5cm of b]          (a) {Auth\\JWT + Google OAuth 2.0\\email verification};
    \node[sub, fill=orange!20, right=0.5cm of a]          (f) {Frontend\\React + Vite\\SSE card-by-card};
  \end{tikzpicture}
};

\end{tikzpicture}
}% end scalebox
\caption{UML Component Diagram: four-layer architecture of ExamGen.
Data flows top-to-bottom during parsing and bottom-to-top during generation.}
\label{fig:arch}
\end{figure}

\subsection{Multi-tenancy and Service Bridge}

\paragraph{Multi-tenancy.}
Each instructor's data is isolated in \texttt{teacher\_\{id\}/} within the dataset
directory. The \texttt{Dataset} class is scoped by \texttt{teacher\_id}, optionally
augmented by a shared corpus (\texttt{include\_shared=True}) when personal archives
are insufficient.

\paragraph{Service bridge (\texttt{docparser\_service.py}).}
This module connects the school management backend (Layer~4) with the generation
engine (Layer~3) by: (i)~launching \texttt{docParser} as an isolated subprocess;
(ii)~converting docParser schema to exam-forge Schema~A; (iii)~writing to the
per-teacher folder; (iv)~invalidating the \texttt{Dataset} cache; and
(v)~triggering cloud RAG indexing.

\subsection{Key Design Choices}

\paragraph{Why Docling?}
PyPDF2 and PDFMiner lack table structure extraction. Docling was chosen for its
TableFormer-based extraction and RTL support, consistent with \citet{mucciaccia2025mcq}'s
preprocessing step that removes noise from PDFs before generation.

\paragraph{Why \texttt{instructor} over prompt-engineered JSON?}
Ad-hoc parsing failed on 30\% of prototype responses.
\citet{mucciaccia2025mcq}'s review process and \citet{morris2024aqg}'s pipeline
both require format correctness; \texttt{instructor}'s self-healing loop provides
this guarantee without output-format fine-tuning.

\paragraph{Why RRF over learned rank aggregation?}
Labelled relevance judgements are unavailable in this domain. RRF is parameter-free
and has been shown to outperform supervised fusion on unseen query distributions
\cite{cormack2009rrf}---precisely the condition encountered in a continuously
expanding examination corpus.
